\practicechapter{第 16 章实践：Benchmark 设计、输入、热路径和报告}

本章对应第一册第 16 章“Benchmark 设计：输入、热路径和报告”。这是第一册实践卷的收束章。你要把前面所有章节的材料合成一个最终研究项目：对一个自己编译的二进制，提交可复查的结构报告、工具交叉验证、符号和反汇编解释、benchmark smoke 与限制说明。

\practicesection{一、对应原书问题：实验系统要能长期维护}

一次 benchmark 可以临时跑，benchmark suite 则要长期维护。长期维护需要固定输入、固定输出 schema、保留 raw data、记录环境和构建、区分正确性失败与性能波动。本卷第一版没有做完整性能平台，但已经具备最小骨架：CMake 构建、GTest/GMock 测试、CLI report、section CSV、Google Benchmark smoke 和人工研究报告。

最终项目要把这些部件连起来，而不是各自截图。报告里每个结论都要能追溯到命令、输出文件、源码或测试。

\practicesection{二、从特例推到规律：快得离谱先怀疑实验}

特例一：某次 benchmark 快到几乎为零，可能是编译器删除了工作，也可能是输入太小、batch 不够、计时边界错了。特例二：一次 Release 比 Debug 慢，可能是环境噪声，也可能是优化改变了代码布局、分支或符号。特例三：不同机器上结果差很多，可能是 CPU、频率、内存、内核权限和后台负载不同。

由此推出规律：benchmark 报告先写实验协议，再写数字，再写解释。解释必须能降级：证据强时可以说“在本环境、本输入、本构建下观察到”；证据弱时只能说“需要进一步复核”。

\practicesection{三、实践任务：最终研究项目目录}

建立如下目录结构：

\begin{lstlisting}[language=bash]
reports/final/
results/final/
materials/final/
\end{lstlisting}

最终项目至少包含九类文件。\code{build.md} 写源码、编译命令、编译器版本和构建类型。\code{binary.report} 保存 \code{cpu1ee\_cli} 的文本报告。\code{sections.csv} 保存 \code{cpu1ee\_cli --sections-csv} 输出。\code{readelf-header.txt} 保存 \cmd{readelf -h} 输出。\code{readelf-sections.txt} 保存 \cmd{readelf -S} 输出。\code{symbols.txt} 保存 \cmd{readelf -s} 或 \cmd{nm -C} 输出。\code{disasm.txt} 保存 \cmd{objdump -d} 或 \cmd{objdump -d -Mintel} 输出。\code{benchmark.txt} 保存 benchmark smoke 输出和限制说明。\code{final-report.md} 从这些证据出发写研究结论。

最终报告要按“事实、解释、限制”三段写。事实只能来自命令输出和测试；解释要说明字段之间如何互相支撑；限制要写当前工具不支持什么、当前 benchmark 没测什么、当前环境可能影响什么。

\practicesection{四、验收：第一册实践闭环}

最终项目通过标准如下：

\begin{enumerate}
  \item 能从源码、构建命令、checksum 定位同一个二进制产物。
  \item 能用本卷解析器和 readelf 交叉验证至少三个字段。
  \item 能解释一个函数的符号、section、反汇编和优化影响。
  \item 能运行 GTest/GMock 测试，并说明它们覆盖的合同。
  \item 能运行 benchmark smoke，并清楚说明它没有测哪些成本。
  \item 最终结论不越界：没有证据的地方写限制，不写成事实。
\end{enumerate}

\begin{rubricbox}
第一册实践卷的合格终点不是“代码能跑”，而是读者能拿一个真实构建产物，从字节、格式、工具、指令、ABI、测试和测量七个角度解释它。目录按原书章节组织，是为了让每个正文概念都在同一个工程里完成一次落地。
\end{rubricbox}
